\documentclass[10pt,twocolumn]{book}
\usepackage[utf8]{inputenc}
\usepackage{geometry}
\geometry{a5paper, margin=1.2cm, columnsep=0.8cm}
\usepackage{graphicx}
\usepackage{epigraph}
\usepackage{lettrine}
\usepackage{xcolor}
\usepackage{booktabs}
\usepackage{enumitem}
\usepackage{fontspec}
\setmainfont{EB Garamond}
\usepackage{titlesec}
\titleformat{\chapter}[display]
{\normalfont\LARGE\scshape}{\chaptertitlename\ \thechapter}{20pt}{\Huge}
\titlespacing*{\chapter}{0pt}{-20pt}{20pt}
\usepackage{fancyhdr}
\pagestyle{fancy}
\fancyhead[LO,RE]{\scshape Thepyrgos}
\fancyhead[RO,LE]{\scshape The Violet Gate}
\fancyfoot[C]{\thepage}
\setlength{\parindent}{1em}
\setlength{\parskip}{0.3em}
\usepackage{tcolorbox}
\tcbuselibrary{skins,breakable}
\newtcolorbox{sidequote}{colback=black!5,colframe=black!30,width=\linewidth,arc=2mm,auto outer arc,boxrule=0.5pt,left=0.5cm,right=0.5cm,fontupper=\itshape}

\begin{document}

\title{\Huge\scshape Thepyrgos\\[0.5cm]\Large The Violet Gate}
\author{Fate's Edge Expansion}
\date{}
\maketitle
\thispagestyle{empty}

\clearpage
\thispagestyle{empty}
\epigraph{Every stair is a question. Every answer is a proof. The Ninth and the Witness argue over whose record holds the truth.}{\textit{--- Brother Konstantios of the Unbroken Lamp, Temple of Light, transcribed 876 AR}}

\tableofcontents
\clearpage

\chapter{The City of Stairs}
\section{Overview}

Thepyrgos is both city and province, a vertical labyrinth of stair-rights, rope-guilds, and synod edicts that shift with the tide. The capital rises in tiers of towers, terraces, and endless stairs above the Dolmis shore---but the stairs are not stone. They are fossilized vertebrae, the spine of a leviathan that Khemesh dragged under when the last Astroegro emperor broke a covenant older than the sea.

Once a Utaran province, now a nation unto itself, Thepyrgos is famed for its universities, libraries, and debating synods, where precedent and philosophy weigh as much as pikes or coin. Yet beneath the scholar's gown and the notary's seal, an older war is fought---between the Witness, who would have every truth recorded and remembered, and the Ninth, who whispers that some truths should never have been written at all.

\begin{sidequote}
The Synod Hall at dusk. Gold-glass mosaics glitter as a Matriarch reads a scroll aloud. \textit{The Beacon Crown sees traitors in the smoke. A child's talent was surrendered to the Commissary last night. Find the informer---or the next stair seized will be yours. And if you hear a stair creak where none exists, do not turn around. The Ninth is counting.}
\end{sidequote}

\section{Geography}

Thepyrgos occupies a peninsula jutting into the Dolmis Sea, bounded by the Black River to the north and the Belmont Marshes to the west. Viterra lies across the Black River, the two realms separated more by law than by water. The countryside is rugged: limestone hills, olive groves, whitewashed villages, and valleys choked with reed-filled marshes. The weather is warmer and wetter than the northern reaches---summers are humid, winters cool and drizzly, and fog rolls in from the sea without warning.

The Astroegro Straits, the narrow sea passage between the Dolmis and the Amaranthine, are controlled by Thepyrgos, making the city a crucial link in the riverine trade between the fractured sun-kingdom and the imperial remnant. A few small colonies exist on the western coast of Hyaar---exiled scholars and retired Chain-Lanterns who trade olive oil and wine for eastern silks and spices. They are said to be even more paranoid than the mother city.

\subsection{The Hyro Peoples}

The Hyro are the surviving descendants of the Astroegro, a human civilization of formidable magic that fell before the Utaran legions ever marched east. Their ruins litter the hills and swamps south of the city, and their legacy haunts every threshold.

A Hyro child born in the hills might dream in geometries that do not exist. A farmer might hear the stones whispering in a language no living throat can shape. The Chain-Lanterns and other Witch-Hunter Commissariats ruthlessly patrol these villages, using posted warrants, candle-smoke tests, and stair-seizures to identify and conscript children with arcane talent. The Hyro have learned to offer substitutes---a memory, a false name, a lock of hair---at crossroads, hoping the Hollow accepts the tribute.

\subsection{The Fall of the Astroegro}

The folk of Thepyrgos hold that Khemesh, the Abyssal Maw, destroyed the Astroegro with a single cataclysmic wave, dragging the emperor's palace into the abyss as punishment for the breaking of an ancient covenant. The scholars of the universities dismiss this as superstition. They point to evidence of slow decline, internal collapse, and the empire's own hubris. But the folk are not far from the truth.

\section{The City}

Thepyrgos is not one city but three, stacked on top of each other like a geological column, each layer pressed upon the one below like sediment in an ocean that never existed. The streets do not run straight. They climb, they curl, they terminate in dead ends that were once doorways to buildings that have since migrated elsewhere. A traveler who walks the same route on consecutive days will find the journey shorter or longer depending on the city's mood. The Thepyrgosi do not speak of this. They simply allow extra time.

\subsection{The Lower City}

The Lower City clings to the harbour like barnacles to a wreck. Its buildings are ancient, their foundations sunk into mud that remembers the weight of Astroegro wharves. The architecture is a palimpsest: Utaran brick atop Astroegro stone atop something older that the Utarans themselves could not identify. Doorways are arched and narrow, designed for bodies smaller than modern measure. Windows are shuttered or bricked entirely. The streets are canals of shadow between leaning walls that have forgotten how to stand straight.

The air smells of brine, cheap wine, and the copper tang of old blood. The fog rolls in from the Dolmis without warning, and when it does, the Lower City becomes a labyrinth of wet stone and muffled footsteps. The fog has a voice, the old folk say. It whispers names. The names are always those of people who died alone.

The rope-guild halls are long, low buildings near the water, their facades carved with scenes of sailors drowning and being rescued and drowning again---the same figures in an endless cycle. Inside, the guilds keep their ledgers, their knots, and their secrets. The ropes they make are not merely rope; they are witnesses. A rope from the Lower City remembers who pulled it, who was tied by it, who was lowered into a grave with it coiled at their feet.

The Commissariat's interrogation cells are cut into the bedrock beneath the harbour master's office. They are reached by a stair that descends nine flights. The ninth flight is below the water table. The walls weep. The prisoners do not know whether the wetness is condensation or the sea trying to claim them early. The Chain-Lanterns do not correct the ambiguity.

\subsection{The Middle City}

The Middle City rises on fossil-stone terraces---the vertebrae of the leviathan that Khemesh dragged under when the last Astroegro emperor broke the covenant. The stone is gray-green, veined with white, and warm to the touch even in winter. In some places, the vertebrae are intact, forming natural arches that the Utarans incorporated into their colonnades. In others, the stone has been carved and recarved so many times that the original bone is visible only as a darker grain beneath the marble.

This is the city of scholars, notaries, and those who live by the weight of the written word. The universities cluster here, their buildings linked by bridges that span the gaps between vertebrae. The bridges are narrow, unrailed, and deliberately unsettling. A student who cannot cross a bridge without looking down, the faculty say, has not yet learned to look up.

The synod halls are domed structures of Astroegro construction, their interiors lined with gold-glass mosaics that depict not saints or emperors but geometric proofs---solved and unsolved. The solved proofs are illuminated in gold leaf; the unsolved are left as dark patches, waiting for a scholar to complete them. The acoustics are enchanted, or perhaps merely a property of the bone-stone. A whisper spoken at one end of the Synod Hall is audible at the other. A lie spoken anywhere in the hall echoes twice---once in the liar's voice, once in truth's. The truth's echo is always softer.

The great archives are windowless, their walls lined with lead to keep out the sea-damp and the Hollow's attention. The shelves are organized not by subject or author but by the date the document was sealed. The oldest documents are at the bottom, in vaults that have not been opened in centuries. The Archivists do not know what those documents contain. They know only that the seals are still intact, and that the lead is warm to the touch.

\subsection{The Upper City}

The Upper City crowns the acropolis, a walled fortress of ancient Astroegro construction. The walls are not stone; they are a single piece of fused material that the Utaran engineers could not replicate and the Thepyrgosi masons cannot repair. When a section cracks, the city does not patch it. It builds a new wall behind the old one, and the old wall is left to lean and whisper. The cracks have a geometry. No one has mapped it.

The Archon's citadel sits at the highest point, a black mass of angles that do not quite meet. The citadel has no windows on its western face. The reason is not recorded. The Archon's chambers are in the eastern wing, where the morning light falls on a mosaic of the Astroegro emperor Ionius the Ninth, his face worn smooth by centuries of dust and deliberate erasure. The Archon does not sit facing the mosaic. No Archon has sat facing it since the fall.

The Beacon Crown's signal towers ring the acropolis wall, their braziers fed with juniper and saltpeter. The smoke from the beacons does not rise; it hangs, forming a permanent gray pall that the Thepyrgosi call the Crown's Shadow. The shadow is thin in summer, thick in winter, and on certain nights---when the wind shifts and the bells are silent---it descends into the Middle City, carrying the taste of ash and the sound of voices that have been dead for generations.

The stairs to the Upper City are guarded by the Palikars, the citizen-militia who answer only to the Synod. The Palikars wear black wool and carry staves of fossil-bone. They do not speak to those they escort. They do not need to. The stairs speak for them.

\subsection{The Stairs}

Between the three cities, the stairs. A thousand stairs, some public, some private, some belonging to guilds or families or no one at all. The Thepyrgosi do not count the stairs. They know that the number changes, and they have learned that a changed number is a warning, not an arithmetic problem.

Each stair has a right---a charter, a tradition, a memory---and to climb without right is to invite the attention of the stair itself. Public stairs may be used by anyone, but they are slow, crowded, and patrolled by the Palikars. Private stairs are faster, but they require a token, a password, or a bribe. Family stairs are the fastest of all, but they are also the most dangerous, for a family stair remembers every member who has climbed it, and it knows when a stranger is pretending.

Some stairs belong to no one. These are the unclaimed stairs, the ones that appear in the gaps between buildings, leading from nowhere to nowhere. The Thepyrgosi avoid them. A stair with no owner, they say, is a stair that the Hollow has already claimed. To climb an unclaimed stair is to arrive at a destination you did not choose.

The stairs have long memories. They remember the Astroegro architects who carved them, working in total darkness because light would have revealed the geometry of the fossil-stone and driven them mad. They remember the Utaran legions who marched up them, counting their steps in Utaran, and the legions who marched down, counting their dead. They remember the fall---the night the leviathan stirred and Khemesh's wave swept the Lower City, drowning a district and leaving behind a new stair that had not been there before.

And sometimes, at night, when the fog is thick and the bells are silent, the stairs remember a step that was never there before. A step that leads down when it should lead up. A step that is warm when it should be cold. A step that, if you place your foot upon it, will take you somewhere that is not Thepyrgos. The Thepyrgosi do not speak of these stairs. They do not avoid them either. They simply pretend not to see them, and the stairs, for their part, pretend not to exist.

\subsection{The Architecture of Dread}

The buildings of Thepyrgos are not designed for human comfort. They are designed for permanence, and permanence, in Thepyrgos, is a kind of threat. Walls are thick enough to muffle screams. Ceilings are high enough to make a whisper echo. Doors are low, forcing even the tallest visitor to bow---a gesture of submission that the building does not acknowledge but nevertheless expects.

The stone is always cold except where it is warm. The warm patches are not near hearths or braziers. They are in corners, in thresholds, in the spaces behind locked doors that have not been opened in living memory. The warmth is not geothermal. The warmth is the memory of something that once lived in the stone and has not yet left.

The windows are narrow, set deep into the walls, and glazed with glass that distorts the view. A person looking out sees the street as a wavering line, the sky as a smear of gray. The distortion is not accidental. The Thepyrgosi do not want a clear view of the outside. A clear view would remind them that the city ends, and what lies beyond the walls---the marshes, the sea, the open road---is less certain than the stone.

The doors are banded with iron, the iron stamped with seals that no longer correspond to any living authority. The seals are still enforced. A door sealed by a dead notary will not open for a living clerk, and the notary's ghost does not accept bribes. The Thepyrgosi have learned to negotiate with ghosts. It is slower than bribery, but the terms are clearer.

\subsection{The City's Mood}

Thepyrgos has a mood, and the mood varies by district. The Lower City is sullen, brooding, its alleys narrow and its silences heavy. The Middle City is anxious, its scholars always looking over their shoulders, its bells always warning of something that has not yet arrived. The Upper City is serene, but it is the serenity of a tomb---a place where nothing changes because nothing is alive enough to change.

At night, the mood shifts. The Lower City becomes watchful, its fog carrying whispers that are not quite words. The Middle City becomes restless, its bells ringing at odd hours for reasons no one can explain. The Upper City becomes silent. Not the silence of sleep, but the silence of a held breath. The Archon's citadel glows faintly, its black walls reflecting starlight that falls from a sky that is never entirely clear.

The Thepyrgosi do not speak of the city's mood. To name it would be to admit that the city has a will, and that will is not entirely benevolent. Instead, they speak of the weather, of the tides, of the bells. They speak of everything except what the city wants from them. And the city, patient as stone, waits for them to ask.

\subsection{The Lethai-thora}

High elf refugees from the fall of the Valewood's Empire of Leaves, the Lethai-thora fled across the Black River and were granted sanctuary in Thepyrgos. Generations later, their descendants still walk the city's tiers, their silver hair and ink-stained fingers a quiet reminder of a pact older than the synod.

The Lethai-thora are the keepers of the deepest archives and the most secret keys. They are not universally trusted---their loyalty is to knowledge, not to Thepyrgos---but they are indispensable. They maintain the forbidden shelves, the sealed vaults, and the catalogs that no one else can read. The institutions of Thepyrgos are human-led, but the Lethai-thora hold a power that transcends mere office: they remember what others have forgotten.

\section{The Universities}
Thepyrgos is home to six universities, each with its own specialisation, its own library, and its own patron. The faculties are human and Lethai-thora, though instructors of other cultures may be found among the more exotic disciplines.

\subsection{The University of Sacred Architecture (Artifice School)}

Geomancy, engineering, bridge-law. The True Masons maintain a chapterhouse here, training engineers who can read the stress-lines of stone and the breaking points of contracts. \textit{Measure twice. Cut once. The mountain remembers the third cut.}

The University of Sacred Architecture occupies a cluster of buildings near the base of the Upper City, where the fossil-stone stairs are oldest and the weight of the citadel presses down. Its lecture halls are vaulted, windowless, and echo with the tap of measuring rods against slate. Students learn to calculate load tolerances in their heads, to sight a foundation's level from a hundred paces, and to hear the difference between a keystone that will hold and one that will whisper for a century before it fails.

The faculty are drawn from the True Masons' guild, and the curriculum is as much oath as instruction. A graduate of Sacred Architecture does not merely know how to build; they know how to seal a bridge against fair-weather promises, how to carve a lintel that remembers the name of every contractor who worked on it, and how to read the stress-lines of a treaty by the way the clauses bear weight. The Weave respects architecture. So does the Hollow. A properly proportioned arch can outlast any curse---but a poorly sealed threshold invites attention.

The school maintains the \textbf{Registry of Keystone Charters}, a brass-bound ledger recording every bridge, wall, and gate sealed by a True Mason's mark. To build without a charter is to invite collapse. To build with a false charter is to invite the Masons' attention. The Registry is kept in a vault behind a door that has no lock---only a riddle. The riddle changes each season. The answer is always a measurement.

\subsection{The Lyceum of Bells and Beacons (Thepyrgosi School)}

Acoustics, signal-law, weather-reading. This school trains the bell-ringers and beacon-keepers who maintain the city's warning systems. Its graduates can read a storm's approach by the taste of the wind and send a message across the province in the time it takes to ring a bell eight times.

The Lyceum is housed in the Beacon Crown itself, a ring of watchtowers on the acropolis where the wind never stops and the bells speak in overlapping dialects. Students learn the grammar of chimes: three short for muster, two long for all clear, a single sustained note for \textit{the Archon is coming, look busy}. They learn to read smoke signals by the way the wind bends them, to interpret mirror-flashes from the harbour, and to distinguish a fog-bell's warning from a fog-bell's lie.

The faculty are retired beacon-keepers, deaf in one ear and patient in both. They teach that sound is a kind of law: it travels, it binds, it cannot be recalled. A bell rung in error can start a war. A bell not rung can lose a fleet. The Lyceum's motto is carved into the largest bell in the Crown: \textit{The note that is not struck is the only note you can take back.}

Graduates of the Lyceum carry a \textbf{Watcher's Bell}, a small bronze chime tuned to the frequency of the Crown's signal network. Once per session, they may ring the bell to send a single word to any other Watcher's Bell within the province. The word arrives in the time it takes the bell to stop ringing. The Weave does not eavesdrop on this communication---but the Hollow listens to bells.

\subsection{The Scriptorium of Sealed Truths (Lethai-thora School)}

Archival science, cryptography, witness-law. The deepest archive in Thepyrgos, where forbidden texts are stored in lead-lined boxes and watched by blind sentinels. The Lethai-thora who work here are divided between those who believe all truths should be preserved and those who believe some truths should never have been written.

The Scriptorium is not marked on any public map. It occupies a warren of rooms behind the Library of Keys, accessible only through a door that requires three keys held by three different Archivists. The building is older than the University district; some say it predates the Utaran occupation, that its walls are Astroegro work, that the lead lining was not added by the Synod but uncovered by it.

Students at the Scriptorium learn the art of the \textbf{sealed statement}: a document that cannot be opened without breaking its own seal, and that will burn its contents if opened by the wrong hand. They learn to read ciphers that change with the tide, to write in inks that vanish when spoken aloud, and to witness oaths without recording them---to hold a truth in memory only, because paper can be stolen and a mind, properly guarded, cannot.

The Scriptorium's divided loyalties are an open secret. The Archivists who serve the Witness believe that every sealed truth should eventually be unsealed, that the world deserves to know what it has forgotten. Those who serve the Ninth believe that some truths were sealed for a reason---that the Astroegro knew what they were doing when they locked their vaults and threw away the keys. The two factions do not speak to each other. They communicate through sealed memoranda. The memoranda are never opened.

Graduates of the Scriptorium receive a \textbf{Sealer's Ring}, a silver band that stamps a wax seal only the wearer can break. Once per session, the wearer may seal a document or a spoken oath with the ring. Any attempt to break the seal---or to violate the oath---triggers a minor backlash (1 Fatigue, the Weave's disapproval) on the violator. The ring does not protect against the Hollow. The Hollow has its own seals.

\subsection{The College of Waters and Weights (Commercial School)}

Alchemy, maritime law, trade logistics. This school trains the factors, notaries, and customs inspectors who keep Thepyrgos' trade routes flowing. Its graduates can calculate interest in their heads and spot a forged seal from across a room.

The College is located near the Salt Gate, where the harbour's tide licks the foundation stones and the air smells of brine and burnt ledgers. Its lecture halls are former warehouses; its practice yards are the docks themselves. Students learn to weigh a cargo by eye, to assess a ship's seaworthiness by the way it sits in the water, and to read a contract's hidden clauses by the spacing of the signatures.

The faculty are retired merchants, former customs inspectors, and a few notaries who have seen enough forgeries to teach the difference between a genuine seal and a clever fake. They teach that the Weave is not the only system of exchange; coin and contract are older than magic, and the Hollow collects debts in all currencies. The College's motto is painted on the wall of every classroom: \textit{The scales do not lie. The people who read them do.}

Graduates of the College may take the \textbf{Merchant's Seal}, a stamp that certifies their authority to inspect cargo, verify weights, and resolve disputes at the Salt Gate. Once per session, a character with the Merchant's Seal may demand to inspect a cargo or a contract without being refused---the Seal carries the weight of the Synod's commercial authority. However, each use ticks the \textbf{Customs' Attention} timer. When the timer fills, the College demands an accounting: a share of the cargo, a favor, or a truthful answer about where the goods came from.

\subsection{The Hyro College of Unfinished Histories (Unofficial, Underground)}

Astroegro lore, folk magic, resistance theory. This college has no buildings, no faculty, and no official recognition. It is a network of Hyro scholars, runaway children, and sympathetic archivists who meet in basements and root cellars to preserve what the Synod would destroy. They teach the old songs, the forgotten proofs, and the names of the Astroegro vaults.

The Hyro College has no fixed location. Meetings are announced by a chalk mark on a stair---a spiral, a triangle, an eye. The mark fades after a day. Those who recognize it gather at dusk in a cellar, a cistern, or the back room of a printer's shop. They bring food, stories, and questions. They leave with answers, warnings, and sometimes a child who needs to disappear.

The curriculum is oral. There are no textbooks because textbooks can be seized. The College teaches the old songs---lullabies in Astroegro that still resonate with the Weave, work chants that keep a well from running dry, mourning verses that lay a ghost to rest without a priest. They teach folk magic: how to tie a knot that holds a promise, how to read a stone's temperature with the back of the hand, how to tell whether a stranger is Hollow-touched by the way their shadow falls.

They also teach resistance. The College is the main conduit for smuggling gifted children out of Thepyrgos to Hyaar or the Vilikari marches. Its members know the hidden stairs, the corruptible guards, the Chain-Lantern patrol schedules. They do not teach violence. They teach escape.

The College has no leader, but it has a symbol: the \textbf{Unfinished Knot}, a length of rope tied in eight loops with the ninth left loose. Those who carry an Unfinished Knot may recognize other members on sight. They may also, once per session, declare a hidden passage or a safe house in any district of Thepyrgos. The GM decides whether the passage is still safe. The Hollow does not know the College's secrets---but it is learning.

\subsection{The Schola of the Unnamed Ninth (Proscribed, Underground)}

Forbidden geometries, unproven theorems, the cult of forgetting. This school does not exist. Its members meet in the Unfinished Stair, in the spaces between bells, in the gaps of the city's memory. They believe that the Astroegro were right to unmake themselves---and that the only way to save the world from the Hollow is to help them finish the job.

The Schola is not a school in any sense the Synod would recognize. It has no faculty, no curriculum, no entrance examination. It has only a text---the \textbf{Ninth Proof}, a fragment of Astroegro geometry that describes how to unmake a single fact without unmaking the facts around it. The Proof is incomplete. The Schola believes that completing it would allow a skilled caster to erase the Hollow from existence by erasing the memory of every debt it has ever collected.

The Schola's members are recruited in dreams, in the margins of forbidden texts, in the silence after a catastrophic backlash. A Threadweaver who fails a dangerous spell badly enough may wake with a single sentence written on their arm in a hand they do not recognize: \textit{The Ninth has seen you. The Stair is open.}

The Schola teaches that the Weave is a lie---that reality is a contract written by powers who cannot be trusted, and that forgetting is not a loss but a liberation. To forget a debt is not to default; it is to refuse to play the game. To forget a name is not to diminish the named; it is to free them from the ledger.

The Synod has outlawed the Schola nine times. The Chain-Lanterns have executed suspected members. The Schola continues because it is not an organisation; it is an idea. And ideas, the Schola teaches, cannot be sealed. They can only be forgotten. And forgetting is their purpose.

There is no mechanical benefit to joining the Schola. There is only a question: \textit{What would you forget, if you could?} The answer is the only tuition. The Weave will remember what you forget. The Hollow will not.

\subsection{The University of Mysteries (The University)}

Magic, anti-magic, threadweaving, metaphysics. This is the oldest and most coveted of Thepyrgos' schools, the one to which every provincial mage aspires and from which every provincial mage is eventually expelled. It trains sanctioned magi and invokers in the deeper mysteries of the world---not merely how to cast, but what casting does to the caster and to the fabric of reality itself.

The University of Mysteries occupies no single building. Its classrooms are scattered across the Middle City: a lecture hall in a converted cistern, a practice yard behind a tannery, a scriptorium whose windows are painted black. The faculty believe that magic should not be comfortable. Comfort leads to carelessness, and carelessness leads to the attention of things that should not be woken.

The curriculum is rigorous and deliberately disorienting. First-year students learn the Three Prohibitions: no blood magic, no summoning without a leash, no working that seeks to unmake a truth the Synod has sealed. Second-year students learn the Three Exceptions: when the Prohibitions may be broken, and who must witness the breaking. Third-year students learn that the Exceptions are lies, that the Prohibitions are suggestions, and that the only real law is the one the caster is willing to enforce with their own life.

The higher you climb, the older and hungrier the world seems. The fourth-year curriculum is not taught; it is discovered. Those who reach the upper tiers find themselves in rooms that were not there the day before, consulting texts that were not written by human hands, and answering questions posed by faculty who may not be entirely alive. Some never come down. The University does not mourn them. It merely adds their names to the register of the disappeared---a ledger kept in a locked room that no one admits exists.

The University is human-led, as are all Thepyrgos institutions, but the Lethai-thora are prominent among its senior faculty. They teach the older courses---the ones that require memory beyond a human lifespan---and they are the ones who decide which students are ready to climb higher. The current Chancellor, a Lethai-thora named Aerendil Thoraen, has held the position for nine decades. He has not aged visibly in that time. He does not speak of his predecessor, and the archives contain no record of an appointment.

\noindent\textbf{Lysandra of the Amber Gate.} No discussion of the University of Mysteries would be complete without naming its most famous---and most controversial---living alumna. Lysandra is a Master of Spellcraft, a Threadweaver of considerable power, and the author of a text the Synod has tried and failed to suppress three times. Her \textit{Threadweaver's Spellbook} is required reading in the University's secret curriculum and forbidden reading in every official syllabus. It teaches the raw grammar of the Weave: intention, tag, and consequence. No patron. No writ. No seal.

Lysandra is not a Runekeeper. She is a free caster---unlicensed, unbound, and unforgettable. She has taught at the University for forty years, despite the Synod's periodic efforts to dismiss her. The Chain-Lanterns have investigated her nine times. Each investigation closed when the lead inquisitor received a letter containing a single sentence: \textit{"The Weave remembers. So do I."}

Her classroom is the Amber Gate, a lecture hall in the lower tiers that smells of old smoke and burnt ambition. She does not assign textbooks. She assigns failure. "Theory is a map," she tells her students. "Practice is the road. You will singe your sleeve. Sew the burnt fabric into a book. One day that book will be the size of a pillow, and you will look at it and know exactly how much you owe."

She is known for her three questions---the questions every Threadweaver must answer before every casting:
\begin{enumerate}[noitemsep]
\item What do you intend? (Be precise.)
\item Which tags serve your intent? (Name them aloud.)
\item What is the risk? (Look at your Position, your Fatigue, your witnesses.)
\end{enumerate}

The University does not officially endorse her methods. The University also does not dismiss her, because the Weave listens when she speaks, and the Weave has a long memory. Her students---the ones who survive---are among the most capable Threadweavers in the province. They are also among the most singed.

The University maintains the \textbf{Registry of Sanctioned Invokers}, a brass-bound ledger that records every licensed magic-user in Thepyrgos. To practice magic without a license is to invite the Chain-Lanterns. To practice with a license is to accept that the University has a claim on your soul---a claim that can be called in at any time. Lysandra's name was struck from the Registry in 852 AR. She still teaches. The Registry still lists her. The contradiction has not been resolved.

Notable faculty include \textbf{Chrysanthos the Unraveler}, a human theorist who believes that all magic is a form of threadweaving and that the Hollow is merely a loose thread waiting to be cut; \textbf{Damek of the Iron Scales}, a former Chain-Lantern who teaches anti-magic and the art of binding rogue casters; and \textbf{Sophia the Silent}, a Lethai-thora archivist who has not spoken in twenty years but who writes answers on water that freeze into legible ice. Students whisper that Sophia is not a teacher but a warning---she climbed too high and saw something that stole her voice.

The University's motto is carved above the door of every lecture hall: \textit{Prove what you can. Seal what you cannot. Forget what should never have been known.} The third clause is the most important. The University does not teach forgetting. It teaches that forgetting is sometimes the only responsible choice.

\noindent\textbf{Mechanical Hook.} A character who trains at the University of Mysteries may gain access to the \textbf{Sanctioned Invoker's Seal}, a brass token that grants lawful permission to practice magic in Thepyrgos. The seal also allows the Chain-Lanterns to track the bearer's location. Removing it is illegal. Keeping it is a leash.

Once per session, a character with a Sanctioned Invoker's Seal may gain +1 die on any sanctioned magic roll---the University's authority shields them from scrutiny. However, each use ticks the \textbf{Registry's Attention} timer. When the timer fills, the University calls in a favor: a task, a testimony, or a truth the character would rather keep hidden.

Characters who study under Lysandra---or who possess a copy of her \textit{Spellbook}---may learn to free cast without the Seal. They gain access to the \textbf{Threadweaver's Grammar}, using the tag-based casting system described in her text. They also lose the protection of the Registry. The Chain-Lanterns will take an interest. The Weave does not care.

\subsection{The Hyro College of Unfinished Histories (Unofficial)}

Astroegro lore, folk magic, resistance theory. This college has no buildings, no faculty, and no official recognition. It is a network of Hyro scholars, runaway children, and sympathetic archivists who meet in basements and root cellars to preserve what the Synod would destroy. They teach the old songs, the forgotten proofs, and the names of the Astroegro vaults.

\subsection{The Schola of the Unnamed Ninth (Proscribed)}

Forbidden geometries, unproven theorems, the cult of forgetting. This school does not exist. Its members meet in the Unfinished Stair, in the spaces between bells, in the gaps of the city's memory. They believe that the Astroegro were right to unmake themselves---and that the only way to save the world from the Hollow is to help them finish the job.

\section{The Unfinished Stair}

\subsection{A Rumor of Ascent}

There is a stair in the University of Mysteries that does not appear on any map. The senior faculty do not speak of it. The junior faculty claim it does not exist. The students know better. They have seen it---or they have seen something---late at night, when the bells have rung the eighth hour and the ninth is still a breath away.

The Unfinished Stair manifests in different places at different times. Some say it emerges from the wall of the Amber Gate lecture hall, where Lysandra's chalk dust hangs in the air like smoke. Others claim it rises from the cistern beneath the Scriptorium, its first step wet, its second dry, its third warm. A few insist they have climbed it from a broom closet in the College of Waters and Weights, emerging from a door that had always been locked and was, upon their return, locked again.

The stair is made of fossil-stone, like so much of Thepyrgos, but the vertebrae are not aligned. They twist, they overlap, they descend into the stone only to rise again a step later. The treads are worn smooth not by feet but by something else---a kind of patient erosion that has nothing to do with friction. The risers are carved with geometries that hurt to look at, angles that should not meet meeting anyway, lines that continue past the edge of the stone into empty air.

You can climb the Unfinished Stair for hours, for days, for what feels like years. The steps do not repeat, but they do not progress either. You will pass the same cracked tread, the same worn carving, the same smear of rust-colored stain. You will tell yourself that you are imagining it. The stair will not correct you.

The rumor---and it is only a rumor, because no one who has climbed to the top has ever returned to describe it---is that the Unfinished Stair has no top. It continues upward forever, through the fossil-stone, through the acropolis, through the sky itself, into a space where the Weave frays and the Hollow's attention sharpens. Those who climb far enough, the rumor says, cease to be themselves. They become part of the stair. Their faces appear in the carvings. Their footsteps echo in the treads. They are not dead. They are not alive. They are pending.

\subsection{Lysandra's Account}

Lysandra of the Amber Gate has climbed the Unfinished Stair three times. She will not speak of the first two ascents. The third, she recorded in a sealed addendum to her \textit{Threadweaver's Spellbook}, which the Synod has attempted to suppress and which her students copy by hand in the hours before dawn.

\begin{tcolorbox}[colback=black!5,colframe=black!30,width=\linewidth,arc=2mm,auto outer arc,boxrule=0.5pt,left=0.5cm,right=0.5cm,fontupper=\itshape,breakable]
\textbf{From the Sealed Addendum, transcribed by a student who shall remain unnamed:}

\textit{The stair found me in the Amber Gate, after the eighth bell, when the fog pressed against the windows and the chalk dust hung still as held breath. I did not look for it. I was erasing a board---a diagram of the Weave's resonance frequencies---and when I turned, the board was gone and the stair was there.}

\textit{I climbed for nine steps before I realized I was climbing. The first step was cold. The fifth step was warm. The ninth step was neither; it was the temperature of skin, and when I placed my foot upon it, I heard my own heartbeat from outside my body.}

\textit{The stair does not measure distance. It measures choices. At each landing, a door. The doors were labeled in a script I could not read, but I understood the labels anyway: REGRET. FORGIVE. FORGET. The door I did not choose was the door that followed me.}

\textit{I climbed for what felt like hours. The fossil-stone vertebrae grew larger as I ascended, or perhaps I grew smaller. The carvings began to move, not quickly, but with the patience of continental drift. A face turned toward me. It was my own face, but older, and its mouth was open as if to speak. It did not speak. It waited.}

\textit{I do not remember reaching the top. I remember reaching a landing where the stair continued in nine directions at once, and the doors were not doors but mirrors, and the mirrors showed not my reflection but the reflections of people who had climbed before me. Some I recognized. Some I did not. All of them were watching.}

\textit{The Weave was thin there. I could see it---not as a metaphor but as a substance, gray and glistening, like the surface of a still pond after rain. And beneath the Weave, something else. Darker. Patient. The Hollow, perhaps, or only the Hollow's shadow.}

\textit{I turned around. The stair behind me was gone. I was standing in the Amber Gate, the chalk dust settling on my shoulders, and the board was clean. The diagram I had been erasing was still visible, but it had changed. The resonance frequencies now included a ninth note. I had not drawn it.}

\textit{I have climbed the Unfinished Stair three times. I will not climb it again. The stair does not need me to climb it. It has my measure. It has everyone's measure. It is only waiting for us to step.}
\end{tcolorbox}

\subsection{Mechanical Notes for the Game Master}

The Unfinished Stair is not a location that can be entered at will. It appears when the Weave is thin---after a catastrophic backlash, during a Convergence of the Ninth, or when a character asks a question that the Synod's archives cannot answer. The GM may trigger a stair manifestation by spending 3+ Story Beats in a scene involving magic, memory, or the Hollow.

When the stair appears, any character who chooses to climb it must make a Spirit + Resolve test (DV 5). On a success, they may ascend one landing and return with a Fragment---a cryptic vision, a forgotten fact, or a single-use Boon. On a failure, they become Lost on the Stair: they vanish for 1d4 scenes, returning with a Condition (Haunted, Marked, or a new Obligation to a Patron of the GM's choice). On a critical failure, they do not return at all. Their name is added to the carvings.

The stair may also be used as a narrative shortcut---a way for characters to bypass a locked door, escape a pursuer, or receive a prophecy. The price is always a memory, a name, or a promise. The Weave does not negotiate. The stair does not forgive. And the Hollow is always, always listening.

\section{The Synod}

The Synod is Thepyrgos' governing body, a council of matriarchs, rectors, and guild-masters who rule by consensus and precedent. Its chambers are in the Synod Hall, a domed building of Astroegro construction where the acoustics are enchanted to carry a whisper to every seat.

The Synod does not have a single leader. The Matriarch presides over sessions but votes only to break ties. Real power lies in the committees: Sacred Measures, Foreign Relations, Unorthodox Texts, and the Commissariat.

The Archon, the city's elected executive, is officially neutral. Unofficially, every Archon in living memory has leaned one way or the other---and each has died under mysterious circumstances before their term ended. The current Archon, Konstantinos, wears a key around his neck. No one knows what it opens. The key is a stair.

\section{The Commissariat}

\subsection{The Chain-Lanterns}

The Chain-Lanterns are the most feared of Thepyrgos' witch-hunting orders. They wear brass lanterns on chains, the lanterns tuned to the frequency of latent magic. A Chain-Lantern patrol can identify a gifted child from a hundred paces by the flicker of candle-smoke.

Their methods are brutal and efficient. The Candle-Smoke Test forces a child to breathe onto a lit candle; if the smoke curls unnaturally, the child is marked. Stair-seizures revoke a family's stair-right, forcing them to use public stairs where the Chain-Lanterns can watch them. Children deemed too dangerous are quieted---their talent surrendered to the Commissary, and they become hollow-eyed servants of the state.

The Chain-Lanterns are not evil. They believe they are saving Thepyrgos from the fate of the Astroegro. They are wrong---but they are not entirely wrong.

\subsection{Other Orders}

The Proctors of the Exam Stair test children for talent in the universities, preferring to recruit rather than suppress. The Vigil of the Unnamed, a secret order within the Commissariat, believes the Chain-Lanterns have gone too far and smuggles gifted children out of the province. The Candle-Keepers maintain the Beacon Crown, watching for signs of heretical magic---and reporting what they see.

\section{The Hyro Villages}

Outside the city, in the limestone hills and reed-choked valleys, the Hyro live much as they have for three thousand years. They farm olives and grapes, keep goats, and fish the coastal lagoons. They speak a dialect of Utaran flavoured with Astroegro words, and they remember---in songs, in stories, in the names they give their children---what the city has chosen to forget.

A Hyro village is built around a well-shrine, a stone circle where offerings are left for the ancestors. The well-shrine is also a hiding place: its deep shaft contains a cache of forbidden texts, a secret passage to the next valley, or a child with the talent who must not be found.

When a Hyro family fears the Chain-Lanterns are coming, they leave an offering at the nearest crossroads: a written name, a lock of hair, a childhood memory. The offering is taken by the Hollow---or by the Ninth, or by something older---and the family is passed over. The child survives, but something is lost: a memory, a name, a future that will never come to pass.

\chapter{The Ninth and the Witness}

\section{The Theological War}

The conflict between the Ninth and the Witness is not a war of swords. It is a war of records---a dispute over whether the world should be remembered or forgotten.

The Witness holds that truth is sacred. Every lie must be exposed, every secret recorded, every broken oath witnessed. The Witness does not forgive; the Witness remembers. Its followers are archivists, investigators, and truth-tellers who believe that the only way to defeat the Hollow is to shine a light into every dark corner. But the Witness has a dark side: its hunger for truth can become cruelty. A Witness fanatic will destroy a family to expose a lie, topple a kingdom to correct the record, and burn a library to ensure that no one reads the wrong book.

The Ninth holds that forgetting is mercy. Some truths are too heavy to carry, some secrets should never have been spoken, and some records are better left unwritten. Its followers are confessors, memory-thieves, and those who have learned that the kindest act is sometimes to forget. But the Ninth has a dark side too: its hunger for forgetting can become annihilation. A Ninth fanatic will burn a library to save a child, erase a name to protect a fugitive, and unmake a person to spare them from a truth they cannot bear.

\subsection{Factions}

The Archivists of the Witness, orthodox and unbending, believe in preservation at all costs. Their leader is the Matriarch of the Ledger, who has never lost a case and never forgiven an error.

The Coterie of the Ninth, radical and secretive, believe that the only way to defeat the Hollow is to forget it exists. Their leader is the Whisper of the Unfinished Stair, a being that may or may not exist, that may or may not be the Ninth itself wearing a mortal face.

The Gray School, a small but influential faction of scholars, argue that the Witness and the Ninth are two halves of a single truth---that remembering and forgetting are both necessary, and that the wise scholar learns when to do which. Their leader is the Silent Rector, who speaks only in writing and has never been seen in public.

\section{The Unfinished Stair}

The Unfinished Stair is the Ninth's domain in Thepyrgos---a hidden flight that appears only to those who have asked a question the Synod cannot answer. The stairs have no end. Each step is a different year. At the top, the Ninth and the Witness are playing knucklebones with forgotten names.

Entering the Unfinished Stair is dangerous not because of monsters, but because of meaning. Each step forces you to confront a truth you have hidden or a memory you have suppressed. Those who climb too high forget why they climbed. Those who climb too low remember everything---and are crushed by the weight.

To find the Unfinished Stair, a character must ask a question that the Synod's archives cannot answer---a question about the Astroegro, about the Hollow, about the nature of truth itself. Then they must walk the city's stairs at midnight, counting their steps. On the ninth flight, the Unfinished Stair will appear---or it will not. The Stair chooses who climbs.

\section{The Astroegro Vaults}

Beneath the city, nine vaults. The Synod knows of three and has sealed them with Witness-branded wax. The Ninth's Coterie has found a fourth. The other five are still buried, waiting for the right key---a true name, a specific death, a memory that has never been spoken.

Each vault contains a weapon of impossible power, a guardian bound to defend it with its last calculation, and a cost---for to open the vault, you must sacrifice something the empire valued. The Ninth seeks to open all nine vaults simultaneously, triggering an unlearning that would erase the concept of record from the world. The Witness seeks to keep the vaults sealed, but remembered perfectly---to write down what happened, even if it must witness the end of everything.

\section{Local Patrons}

The Unwritten Citation is the Witness's aspect in Thepyrgos---a precedent that exists only in the margin, a law that was never written but is still enforced. Its servants are notaries, archivists, and truth-tellers who believe that every lie leaves a stain.

The Silent Stair is the Ninth's aspect---the step that is never counted, the threshold that opens only once. Its servants are confessors, memory-thieves, and those who have learned that the kindest act is sometimes to forget.

The Bell-Founder is Mykkiel, the law that rings true even when it hurts. The Abyssal Pressure is Khemesh, the weight beneath the city that the deep remembers. The Lantern Bearer is the Traveler, the light that guides lost ships and the lost. The Astroegro Remnant is the Sacred Geometry---the proof that was never finished, the equation that still ticks.

\chapter{Mechanics}

\section{Core Timers}

\subsection{Unfinished Stair Timer [6]}

Tracks the party's entanglement with the hidden architecture of Thepyrgos. Ticks when the party uses a hidden stair, asks a question the Synod cannot answer, or speaks a true name of an Astroegro ruin.

At 6 segments, the party discovers a new hidden stair---one that leads to an Astroegro Vault---but they are also marked. The Witness's attention falls on them. Their reflections in mirrors begin to lag by a heartbeat. The Chain-Lanterns receive a vision of their faces.

Reduce the timer by feeding a memory to the Ninth's Coterie (lose a fact, a skill point, or a Bond) or by recording a truth in the Synod's archives (gain the Witness's Annotation).

\subsection{Synod's Patience [4]}

Tracks the Synod's tolerance for the party's meddling. Ticks when the party violates a stair-right, publishes a forbidden text, or interferes with a Chain-Lantern investigation.

At 4 segments, the Synod issues a Writ of Censure. The party loses access to one privilege (library access, safe passage, an ally's aid) until they make amends.

\subsection{Heresy Index [4]}

Tracks the Commissariat's suspicion of the party. Ticks when the party harbours a gifted child, visits a Hyro village after a raid, or uses unsanctioned magic.

At 4 segments, the Chain-Lanterns conduct a Candle-Smoke Test on the party. If they fail, they are marked---gain a Condition (Marked by the Commissariat) that gives -1 die to all social rolls with authorities until cleared.

\section{The Candle-Smoke Test}

When the Chain-Lanterns test a village, roll 1d6:
\begin{itemize}[noitemsep]
\item 1--2: They find a gifted child. The party must intervene or witness a tragedy.
\item 3--4: They find a substitute and leave.
\item 5--6: The village has already paid in silence. No raid occurs.
\end{itemize}

The party may influence the roll by offering their own substitute. Spend a Boon to shift the result by 1.

\section{Hedge Gift: The Offering at the Crossroads}

Once per session, a character who has learned the Hyro customs may leave a substitute at a crossroads---a written name, a lock of hair, a childhood memory. For the next scene, they are invisible to the Chain-Lanterns (cannot be tracked or identified by candle-smoke). The substitute becomes an Unfinished Proof that the GM may call in later. The Hollow accepted it. The Hollow will ask for more.


\chapter{Lore Echoes}

The threads that bind Thepyrgos to the wider world are many, and some of them are old enough to have forgotten their own origins. The following fragments are drawn from the archives of the Scriptorium, the marginalia of Lysandra's \textit{Spellbook}, and the whispered confessions of travelers who have crossed the Black River and returned changed.

\section{The Bitter Root}

Daughters of the Cord run a hospice in the lower tiers, tending to the victims of the Chain-Lanterns' purges. The building is a gray stone structure near the Salt Gate, its door always unlocked, its threshold worn smooth by the feet of the desperate. No one knows how it has survived so long. The Chain-Lanterns have raided it nine times. Each time, they found nothing---no fugitives, no forbidden texts, no evidence of the unlicensed magic they sought. The sisters do not explain. They simply offer tea and ask the inquisitors to wipe their boots.

The hospice specializes in a particular kind of healing. The Daughters do not mend wounds. They cut threads---the invisible cords that bind a person to a debt, a curse, or a Hollowed's attention. A patient who enters the hospice may leave with their name erased from a ledger, their obligation transferred to a sister who will carry it to the grave, or their memory of a traumatic event replaced with a dream of rain. The cost is not paid in silver. It is paid in stories, in silences, and in promises that the patient will one day cut a thread for someone else.

A renegade Sister hides among the bell-runners of the Middle City, her gray robes traded for a tinker's apron and a cap that shadows her face. She is trying to unbind her Hollowed---a child she could not save, whose name she carries in a locket that never warms. She believes that if she can find the original entry in the Commissariat's registry, she can cut the thread that ties the child to the Hollow. She is probably wrong. She is certainly desperate. The bell-runners protect her because she once cut a thread for one of their own.

\textbf{Mechanical Hook:} A character who seeks the Daughters' aid may reduce their Obligation to any Patron by 2 segments, or remove a single Condition caused by a curse or a broken oath. In exchange, they gain a \textbf{Thread to Cut} ---a promise that they will perform a similar service for a stranger before the next full moon. If they do not, the thread tightens, and the Hollow gains a Story Beat to spend on their next scene.

\section{The Threshold Compendium}

Saikou Ira's unfinished masterwork, the \textit{Threshold Compendium}, classifies places where the Weave wears thin and the Hollow's attention sharpens. The Unfinished Stair is listed as a Class IV Anchor site---not merely a threshold but a \textit{persistent instability}, a wound in reality that has not healed because the Astroegro Empire's last calculation is still ticking inside its walls.

According to the \textit{Compendium}, Saikou Ira himself visited the Unfinished Stair during his final expedition. He left a copper coin on every ninth step, a practice he believed would "anchor the ascent" and prevent the stair from migrating further into the city. The coins are still there. Those who have climbed the stair report seeing them---tarnished, warm, and arranged in patterns that shift between ascents. A climber who takes a coin finds that the stair follows them home. It appears in dreams, in the corners of mirrors, in the moment between wakefulness and sleep when the floor seems to tilt.

The \textit{Compendium} also notes that the Unfinished Stair is not the only Class IV Anchor in Thepyrgos. The Beacon Crown's shadow, the Blue Cistern beneath the forum, and the sealed vaults of the Scriptorium all register as Class III or higher. Saikou Ira intended to map them all. He disappeared before completing the work. His body was never found, but his brass compass was recovered from the ninth step of the Unfinished Stair, its needle spinning without rest.

\textbf{Mechanical Hook:} A character who studies the \textit{Threshold Compendium} (or a fragment thereof) gains +1 die on any roll to navigate the Unfinished Stair or to identify a thin place in Thepyrgos. However, the \textit{Compendium} is a cursed text; each time it is consulted, the GM gains 1 Story Beat, as the Hollow notes that someone is learning its geography.

\section{Malachai, the Chained Angel}

A silver coin that never spends circulates among the Hyro villages of the southern hills. It is old---older than the Utaran occupation, older than the Synod's seal. The obverse bears the profile of the last Astroegro emperor, his face worn smooth by centuries of handling. The reverse shows a spiral that seems to turn when the coin is held at the edge of vision. The coin does not clink when dropped. It lands with the soft thud of flesh on stone.

The tradition is this: when the Chain-Lanterns approach a village, a family may offer the coin to the Hollow. They place it on the doorstep, speak the name of their child, and retreat into the house. In the morning, the coin is gone. The Chain-Lanterns pass over that house. The child is safe---for now. But the coin will reappear in the mouth of a drowned sailor, under the tongue of a stillborn goat, or pressed into the palm of a traveler who asks no questions. The coin is the price the Hollow demands for forgetting. And the Hollow always collects interest.

Those who have taken the coin---who have carried it in their purse, spent it on a loaf of bread (the baker will always accept it, and the bread will always taste of ash), or passed it to a stranger in the night---report dreams of the last Astroegro emperor. He is not angry. He is not kind. He is simply waiting, and in his waiting, he asks a single question: \textit{What would you forget, if you could?}

The Hyro do not answer. They pass the coin along.

\textbf{Mechanical Hook:} A character who possesses the silver coin may, once per session, automatically succeed on a single roll to avoid detection by the Chain-Lanterns or to hide from the Hollow's attention. After the roll, the coin vanishes from their possession and reappears elsewhere in the world. The GM gains 3 Story Beats to spend on a complication involving the coin's new owner.

\section{The Gray Wanderer's Grimoire}

Thepyrgos psions are called Stair-Readers. They are not licensed by the Synod, nor are they hunted by the Chain-Lanterns. They exist in a gray space---tolerated because they are useful, feared because they are uncontrollable. A Stair-Reader climbs the Unfinished Stair in their dreams, trading memories for glimpses of the Astroegro vaults, the location of a lost archive, or the name of a child who has not yet been taken.

The Gray Wanderer---a being of considerable power and uncertain origin---climbed the Unfinished Stair once in the waking world. He will not speak of what he saw. He will only show the brass key that hangs from a cord around his neck. The key has no lock. The teeth are worn smooth, as if they have been trying to open something for a very long time. When asked what the key opens, the Gray Wanderer says, \textit{The lock is the question I forgot to ask.}

The Synod has tried to confiscate the \textit{Grimoire} three times. Each time, the inquisitor assigned to the task dreamed of the Unfinished Stair that night and woke with a brass key around their neck and no memory of how it got there.

\textbf{Mechanical Hook:} A character who studies the \textit{Gray Wanderer's Grimoire} may, once per session, reroll a failed roll involving the Unfinished Stair, a thin place, or a dream-omen. On the reroll, any die that shows a 1 is automatically rerolled again (the Wanderer's marginalia corrects the error). However, the character gains a \textbf{Key Without a Lock} ---a brass key that appears in their pocket after the roll. The key does nothing. It simply waits. The GM gains 1 Story Beat per session until the key is passed to another.

\section{The Ninth}

In any Thepyrgos scene, the ninth bell rung in an hour opens a threshold. The wise stop at eight and cover their ears. Those who ring the ninth---by accident, by curiosity, or by the carelessness of a bell-ringer who has forgotten the old warnings---are said to see a pale figure standing on the Beacon Crown's mirror tiles, waiting to receive a name.

The figure is the last Astroegro emperor, Ionius the Ninth. He does not move. He simply waits, and in his waiting, he asks a question that those who see him cannot help but answer: \textit{What is your name?}

Those who answer truthfully find that their name is added to a ledger that never closes. They are not cursed. They are not blessed. They are merely \textit{recorded}. And the Hollow, when it comes to collect, will know exactly where to find them.

Those who do not answer---who turn away, who cover their faces, who run---find that the figure follows them. Not quickly. Not menacingly. Simply \textit{there}, at the edge of vision, in the reflection of a window, in the shadow of a stair. The figure does not demand an answer. It is patient. It has been waiting for three thousand years.

The Thepyrgosi have a saying: \textit{The ninth bell is the Hollow's dinner bell. Ring it if you are hungry. Ring it if you are not afraid to be the meal.}

\textbf{Mechanical Hook:} In any scene set in Thepyrgos, the GM may spend 2 Story Beats to have the ninth bell ring. All characters within the scene must make a Spirit + Resolve test (DV 4). On a failure, they see the pale figure and must either speak their name aloud (gaining a \textbf{Recorded} Condition: the Hollow knows where to find them) or flee in terror (losing their next action). On a success, they shake off the vision---but the GM gains 1 Story Beat as the figure notes their resistance.

\chapter{Superstitions}

These superstitions are observed by the Thepyrgosi, the Hyro, and the more cautious visitors. They may be invoked once per session. Burn the listed resource to gain +1 die on the specified roll.

\noindent\textbf{Bread to the Stair:} Toss a crust onto a rung before climbing. Ignore the first Tremor penalty.

\textbf{Knot the Rope:} Tie a single knot in a crane line. Cancel Rope Guild Interdict or take --1 die on your next lift (your choice).

\textbf{Name to the Smoke:} Whisper a hunter's name into the Beacon Crown's fire. Gain +1 Effect when evading the Chain-Lanterns this scene, but gain the Synod's Notice.

\textbf{Offer the Substitute:} Leave a written false name at a crossroads. The next Candle-Smoke Test against you automatically fails---but the Hollow gains an Unfinished Proof on you.

\textbf{Hold the Eighth Breath:} Before crossing any threshold in Thepyrgos, hold your breath for eight heartbeats. The ninth heartbeat is the threshold's. If you breathe on the ninth, the threshold will remember you. (Mechanically: gain +1 die on the next roll to avoid being tracked through a doorway or gate. On a miss, the GM gains 1 Story Beat.)

\textbf{Leave a Copper for the Stair:} Place a copper coin on the ninth step of any stair you descend. The coin will be gone by morning. The stair will not follow you home. (Mechanically: after descending any significant stair in Thepyrgos, spend 1 Boon to avoid being followed or tracked for the remainder of the scene.)

\vfill
\begin{center}
\textit{The road is south. The ink is fresh. The lamp is borrowed.}\\
\textit{In Thepyrgos, the stairs remember. The bells witness.}\\
\textit{And the Ninth is always counting.}
\end{center}

\end{document}
